\chapter{Poker Without Revealing the Cards}

\label{chap:poker-without-revealing-the-cards}

\researchstatus{Educational toy construction}{The present C program runs in one
process and illustrates a staged design.  It does not establish private
distributed execution, a verifiable shuffle, malicious security, fairness, or
production readiness.}

\begin{abstract}
This chapter studies a candidate decomposition of the dealing phase of
dealer-free Texas Hold'em.  The objective is not yet to implement the full
game, but to expose four questions that a complete construction would have to
answer: how the deck is randomized, how several parties influence its order,
how cards are delivered, and how a later opening can be checked against an
earlier digest.  The accompanying program is an instructional, single-process
demonstrator.  It makes the stages inspectable; it does not prove that the
resulting procedure has the confidentiality, integrity, or fairness properties
required of a deployed protocol.
\end{abstract}

\section{Problem Setting}
Suppose a collection of players wishes to emulate the role of a trusted dealer
in Texas Hold'em. A satisfactory protocol would need at least the following
properties:
\begin{enumerate}
    \item \textbf{Card privacy.} No player should learn another player's hole
    cards before the designated reveal phase.
    \item \textbf{Dealing integrity.} No card may be duplicated, omitted, or
    substituted without detection.
    \item \textbf{Transcript consistency.} The participants must agree on a
    single public record of the shuffled deck and the subsequent public cards.
    \item \textbf{Showdown verifiability.} When a player eventually reveals a
    hand, the rest of the table must be able to verify that the hand was fixed
    in advance.
\end{enumerate}

The present chapter restricts attention to a model of the dealing problem.  It
asks how cards might be generated, shuffled, assigned, and later opened,
without yet formalizing betting rounds or final winner determination.  The
current program simulates all parties inside one address space, so its printed
``private'' views are presentation boundaries rather than confidentiality
boundaries.

\section{Assumptions and Scope}
The current demonstrator makes three explicit assumptions.
\begin{enumerate}
    \item Each player has access to local randomness.
    \item The players can publish values to a shared transcript visible to all
    participants.
    \item The implementation uses a commutative-encryption toy model in order
    to expose the protocol structure clearly.
\end{enumerate}

The third point is important. The demonstrator is useful as a protocol sketch
and as an executable reference for the staged construction, but it should not
yet be interpreted as a hardened cryptographic system. In particular, no claim
is presently made regarding resistance to malicious deviations, adaptive
adversaries, collusion strategies, or implementation-level side channels.

\section{Four-Stage Construction}
We organize the candidate construction into four stages so that each
requirement can be
examined independently.
\begin{enumerate}
    \item \textbf{Intended deck randomization.} Form a permutation of a
    canonical 52-card deck and identify the resulting in-memory state.
    \item \textbf{Toy multi-party transformation.} Simulate each player
    transforming the deck so that the intended joint-control structure is
    visible.
    \item \textbf{Toy dealing.} Simulate delivery of hole cards and publication
    of community cards.
    \item \textbf{Illustrative hand digest and opening check.} Record a digest
    of each simulated hand and later recompute it from the revealed inputs.
\end{enumerate}

\section{Stage 1: Intended Deck Randomization}
Let \(D = (c_1,\dots,c_{52})\) denote the canonical deck. The first stage
applies a random permutation to \(D\), thereby producing a deck order that is
not publicly revealed. A public digest of the resulting state is then
printed.  Within this execution, the digest serves as a checkpoint derived
from the array state used by later stages.  The program does not provide a
network transcript, and the
short non-cryptographic deck digest printed by the demo is not evidence that a
remote party is bound to a hidden permutation.

\section{Stage 2: Toy Multi-Party Transformation}
Each player next applies two transformations to the entire deck: a private
shuffle and a private commutative-encryption layer. Because the encryption
operations commute, the layers may later be removed in different orders
without changing the final plaintext card in the toy arithmetic.  Both players
and the deck nevertheless live inside one process.  Consequently, this stage
illustrates the intended order of operations; it does not implement a private
or verifiable multi-party shuffle.

After each simulated player finishes the transformation, the program prints a
short digest of the encrypted deck.  These values are useful execution
checkpoints.  They are not presented here as hiding or binding commitments
against an adversarial participant.

At the end of this stage, the program also prints a cryptographic digest for
each encrypted card that is about to be dealt.  Later recomputation checks
whether the encoded in-memory sample produces the recorded digest during an
honest execution.  A
complete protocol would still need canonical messages, session binding,
participant authentication, and a security argument before these digests could
support a claim about substitution by a malicious party.

\section{Stage 3: Toy Dealing}
The dealing phase consumes cards from the jointly shuffled encrypted deck.
Private hole cards are obtained by removing encryption layers in an order that
would leave the final plaintext visible only to the intended recipient in a
distributed realization.  In the current program, one process computes and can
access every result.  Community cards are obtained by removing all toy layers
and are then printed.

At the end of this phase, a future transcript should support checks of the
deal's internal consistency. In the current demonstrator this is checked by verifying
that no plaintext card identifier appears twice among the dealt hole cards and
community cards. In addition, each drawn encrypted card is checked against the
digest printed in Stage~2.  These are implementation invariants for one run,
not a proof of dealing integrity under adversarial execution.

\section{Stage 4: Illustrative Hand Digest and Opening Check}
Once the hole cards have been assigned in the simulation, each player \(P_i\) samples a
private nonce \(N_i\) and publishes a digest
\[
    C_i = H\!\left(N_i \parallel 2 \parallel h_{i,1} \parallel h_{i,2}\right),
\]
where \(h_{i,1}\) and \(h_{i,2}\) are the player's hole cards and \(H\) is a
cryptographic hash function.  In the implementation, the nonce has a fixed
length, the count is one byte, and each card is encoded as a fixed-width
little-endian integer.  The displayed formula is shorthand for that typed,
fixed-length encoding rather than unrestricted string concatenation.

At showdown, the program reveals \(N_i\) together with the simulated private
cards and recomputes \(C_i\).  A matching value shows only that the disclosed
inputs recompute to the earlier digest inside that run.  Calling this a hiding and
binding commitment between distrustful network participants would require a
complete message format, session and identity binding, an adversary model, and
a proof or reduction that this chapter does not yet provide.

\section{What the Decomposition Is Intended to Isolate}
The four-stage decomposition separates four central questions behind dealer-free
poker:
\begin{enumerate}
    \item secrecy of the final deck order,
    \item joint control over the shuffle,
    \item the goal of private delivery of hole cards,
    \item and the goal of checking a delayed opening of a private hand.
\end{enumerate}

The demonstrator lets a reader inspect these questions in sequence.  It does
not establish that the questions have been solved under distrust.  A complete
construction must additionally specify how actions are recorded, how betting
rounds are synchronized, how hand values are computed, and how disputes are
resolved when one or more parties deviate from the stated protocol.

\section{Implementation Companion}
The executable demonstrator associated with this chapter lives under
\begin{center}
\path{src/poker_without_revealing_the_cards/}
\end{center}
Its primary source file is
\begin{center}
\path{poker_without_revealing_the_cards.c}.
\end{center}
It can be run with:
\begin{verbatim}
make -C demostrations demo_poker
\end{verbatim}

The implementation prints stage-by-stage logs, deck digests, encrypted deck
previews, dealt cards, and showdown validation data. Its purpose is to make the
construction inspectable rather than to claim final protocol completeness.

\section{What Cryptography Does Not Solve}
Even a future dealer-free cryptographic construction would not decide whether
the game is fair in a social or economic sense, prevent a player from refusing
to continue, guarantee network availability, establish human identity, or make
collusion irrational.  The present program does less: it does not isolate
players' memories, transport authenticated messages, prove correct shuffles,
or enforce simultaneous progress.  Cryptographic mechanisms can protect
specified information and invariants only after those properties and the
adversary have been defined.

\section{Limitations and Next Steps}
The current chapter and implementation deliberately stop short of a full poker
system. The following components remain open:
\begin{enumerate}
    \item betting and action privacy,
    \item hand-ranking evaluation,
    \item authenticated network transport,
    \item protection against malicious rather than merely cooperative or
    semi-honest execution,
    \item canonical, session-bound protocol messages,
    \item verifiable shuffle correctness,
    \item a precise commitment construction and proof for the chosen model,
    \item fairness, abort handling, and recovery,
    \item and isolation of each participant's private state.
\end{enumerate}

These omissions are not accidental; they define the next research layers. The
current contribution is an inspectable decomposition of the dealing problem
into stages that can be refined and challenged independently.
